%
% frontmatter/abstract.tex
%
\chapter*{Abstract}
\addcontentsline{toc}{chapter}{Abstract}

The need for geographic information extraction from unstructured texts has become crucial to many areas such as urban planning and disaster response. In this thesis, we compare several machine learning methods for automatic geotag recognition (which can be considered an aspect of named entity recognition [NER]) when applied to a large multilingual dataset of complaints about public transportation from citizens of Kazakhstan who speak Kazakh or Russian; this dataset was generated by City Transportation Systems (CTS) in Astana, Kazakhstan. We have implemented four different approaches to recognize the location of each record of complaints in the dataset: (1) A few-shot Large Language Model (LLM) pipeline utilizing Qwen2.5 (7B), which was deployed on-site through Ollama, (2) A few-shot Large Language Model pipeline utilized the same method as above but with a fine-tuned Multilingual-E5-base (Me5-base) encoder, (3) Fine-tuning of XLM-RoBERTa, (4) Fine-tuning of ruBERT.

There were six domain specific types in the dataset: \texttt{ STREET\_NAME }, \texttt{ROUTE\_NUM}, \texttt{BUS\_ID}, \texttt{PLATE\_NUM}, \texttt{STOP\_NAME}, \texttt{STREET\_INTERSECTION}. Each type was tagged according to the BIO tagging scheme. Results from our experiments show that the Qwen2.5 LLM achieved the highest precision (0.788) and macro F1 score ( 0.635); it was also the only model able to identify the complex \texttt{STREET\_INTERSECTION} tag. The Me5-base achieved the best recall (0.653). The XLM-RoBERTa had the greatest ability to identify a particular route number tag (F1 = 0.920), while ruBERT's (F1 = 0.223) lack of success in identifying code switching multilinguality data resulted in its being unable to correctly perform NER in this environment. 


\bigskip
\keywords{Named Entity Recognition, Geotag Extraction, XLM-RoBERTa, BERT, Large Language Models, Multilingual NLP, Kazakh, Russian, Transport Complaints, BIO Tagging, Sequence Labeling}

\clearpage
